% !TeX program = xelatex
\documentclass{article}
\usepackage{kotex}
\usepackage[a4paper,left=3cm, right=3cm, top=2cm, bottom=2cm]{geometry}
\usepackage{amsmath}
\usepackage{graphicx}
\usepackage{caption}
\usepackage{setspace}
\usepackage{xcolor}
\usepackage{titlesec}
\usepackage{amssymb}
\usepackage{tcolorbox}
\usepackage{wrapfig}
\usepackage{amsthm}
\usepackage{multirow}
\usepackage{array}
\usepackage{tikz}
\usepackage{longtable}
\usepackage{pgfplots}
\pgfplotsset{compat=1.18}

\title{1.6 More on Linear Systems and Invertible Matrices}
\date{}
\author{}
\setstretch{1.3}

\titleformat{name=\section, numberless}
  {\normalfont\large\bfseries\color{blue}}
  {}
  {0pt}
  {}
\geometry{a4paper, margin=1in}

% Red subheading style (matches the textbook's red section headers)
\newcommand{\redheading}[1]{{\vspace{2mm}\noindent\Large\bfseries\color{red!55!black} #1\par}\vspace{2mm}}

\newtheoremstyle{mystyle}
  {}{}{\itshape}{}{\bfseries}{.}{.5em}{}
\theoremstyle{mystyle}

\begin{document}
\maketitle

{\itshape\color{blue!45!black}
This section shows how the inverse of a matrix can be used to solve a linear system, and develops further results about invertible matrices.\par}

\redheading{Number of Solutions of a Linear System}

Section 1.1 asserted (from Figures 1.1.1, 1.1.2) that every linear system has zero, one, or infinitely many solutions. We can now prove it.

\begin{tcolorbox}[colback=red!4, colframe=red!60!black, title=Theorem 1.6.1, boxrule=0.5mm, arc=3mm]
A system of linear equations has zero, one, or infinitely many solutions. There are no other possibilities.
\end{tcolorbox}
\textbf{Proof} For $A\mathbf{x}=\mathbf{b}$, exactly one of ``no solutions,'' ``exactly one,'' or ``more than one'' holds; it suffices to show the last case gives infinitely many. Suppose $\mathbf{x}_1,\mathbf{x}_2$ are distinct solutions and let $\mathbf{x}_0=\mathbf{x}_1-\mathbf{x}_2\ne\mathbf{0}$. Then $A\mathbf{x}_0=A\mathbf{x}_1-A\mathbf{x}_2=\mathbf{b}-\mathbf{b}=\mathbf{0}$, so for any scalar $k$,
\[
A(\mathbf{x}_1+k\mathbf{x}_0)=A\mathbf{x}_1+k(A\mathbf{x}_0)=\mathbf{b}+k\mathbf{0}=\mathbf{b}
\]
Since $\mathbf{x}_0\ne\mathbf{0}$ and $k$ ranges over infinitely many scalars, $A\mathbf{x}=\mathbf{b}$ has infinitely many solutions. $\blacksquare$

\redheading{Solving Linear Systems by Matrix Inversion}

Beyond Gauss--Jordan and Gaussian elimination, invertibility gives an explicit \emph{formula} for the solution.

\begin{tcolorbox}[colback=red!4, colframe=red!60!black, title=Theorem 1.6.2, boxrule=0.5mm, arc=3mm]
If $A$ is an invertible $n\times n$ matrix, then for every $n\times1$ matrix $\mathbf{b}$, the system $A\mathbf{x}=\mathbf{b}$ has exactly one solution: $\mathbf{x}=A^{-1}\mathbf{b}$.
\end{tcolorbox}
\textbf{Proof} $A(A^{-1}\mathbf{b})=\mathbf{b}$ shows $\mathbf{x}=A^{-1}\mathbf{b}$ is a solution. If $\mathbf{x}_0$ is any solution, then $A\mathbf{x}_0=\mathbf{b}$, and multiplying by $A^{-1}$ gives $\mathbf{x}_0=A^{-1}\mathbf{b}$ -- so it's the only one. $\blacksquare$

\subsection*{EXAMPLE 1 | Solution of a Linear System Using $A^{-1}$}
Solve $x_1+2x_2+3x_3=5,\ 2x_1+5x_2+3x_3=3,\ x_1+8x_3=17$.

\textbf{SOLUTION:} In matrix form $A\mathbf{x}=\mathbf{b}$ with $A=\begin{bmatrix}1&2&3\\2&5&3\\1&0&8\end{bmatrix}$. From Example 4 of Section 1.5, $A^{-1}=\begin{bmatrix}-40&16&9\\13&-5&-3\\5&-2&-1\end{bmatrix}$, so
\[
\mathbf{x}=A^{-1}\mathbf{b}=\begin{bmatrix}-40&16&9\\13&-5&-3\\5&-2&-1\end{bmatrix}\begin{bmatrix}5\\3\\17\end{bmatrix}=\begin{bmatrix}1\\-1\\2\end{bmatrix}
\]
i.e., $x_1=1,x_2=-1,x_3=2$. (This method only applies when the system is square and the coefficient matrix is invertible.)

\redheading{Linear Systems with a Common Coefficient Matrix}

For a sequence $A\mathbf{x}=\mathbf{b}_1,\dots,A\mathbf{x}=\mathbf{b}_k$ sharing an invertible $A$, the solutions $\mathbf{x}_i=A^{-1}\mathbf{b}_i$ need one inversion and $k$ multiplications. More efficiently, form
\[
[A\mid\mathbf{b}_1\mid\mathbf{b}_2\mid\cdots\mid\mathbf{b}_k] \tag{1}
\]
and reduce to reduced row echelon form by Gauss--Jordan elimination, solving all $k$ systems at once -- and this works even when $A$ is not invertible.

\subsection*{EXAMPLE 2 | Solving Two Linear Systems at Once}
Solve (a) $x_1+2x_2+3x_3=4,\ 2x_1+5x_2+3x_3=5,\ x_1+8x_3=9$ and (b) $x_1+2x_2+3x_3=1,\ 2x_1+5x_2+3x_3=6,\ x_1+8x_3=-6$.

\textbf{SOLUTION:} Both share the coefficient matrix of Example 1. Augmenting with both constant columns:
\[
\left[\begin{array}{ccc|cc}1&2&3&4&1\\2&5&3&5&6\\1&0&8&9&-6\end{array}\right]
\ \xrightarrow{\text{reduce to RREF}}\
\left[\begin{array}{ccc|cc}1&0&0&1&2\\0&1&0&0&1\\0&0&1&1&-1\end{array}\right]
\]
The last two columns give the solutions: (a) $x_1=1,x_2=0,x_3=1$; (b) $x_1=2,x_2=1,x_3=-1$.

\redheading{Properties of Invertible Matrices}

To show $A$ is invertible has meant finding $B$ with $AB=BA=I$. In fact, either equation alone forces the other.

\begin{tcolorbox}[colback=red!4, colframe=red!60!black, title=Theorem 1.6.3, boxrule=0.5mm, arc=3mm]
Let $A$ be square.\\
(a) If $BA=I$ for square $B$, then $B=A^{-1}$.\\
(b) If $AB=I$ for square $B$, then $B=A^{-1}$.
\end{tcolorbox}
\textbf{Proof (a)} It suffices to show $A$ is invertible; then multiplying $BA=I$ on the right by $A^{-1}$ gives $B=A^{-1}$. By Theorem 1.5.3, it suffices to show $A\mathbf{x}=\mathbf{0}$ has only the trivial solution. If $A\mathbf{x}_0=\mathbf{0}$, multiplying on the left by $B$ gives $BA\mathbf{x}_0=B\mathbf{0}$, i.e., $\mathbf{x}_0=\mathbf{0}$. $\blacksquare$

\redheading{Equivalence Theorem}

We can now add two more statements to Theorem 1.5.3.

\begin{tcolorbox}[colback=red!4, colframe=red!60!black, title=Theorem 1.6.4 -- Equivalent Statements, boxrule=0.5mm, arc=3mm]
For an $n\times n$ matrix $A$, the following are equivalent:\\
(a) $A$ is invertible.\quad (b) $A\mathbf{x}=\mathbf{0}$ has only the trivial solution.\\
(c) The reduced row echelon form of $A$ is $I_n$.\quad (d) $A$ is a product of elementary matrices.\\
(e) $A\mathbf{x}=\mathbf{b}$ is consistent for every $n\times1$ matrix $\mathbf{b}$.\\
(f) $A\mathbf{x}=\mathbf{b}$ has exactly one solution for every $n\times1$ matrix $\mathbf{b}$.
\end{tcolorbox}
\textbf{Proof} Since (a)--(d) are already equivalent (Theorem 1.5.3), it suffices to show $(a)\Rightarrow(f)\Rightarrow(e)\Rightarrow(a)$.

$(a)\Rightarrow(f)$: This is Theorem 1.6.2.

$(f)\Rightarrow(e)$: Immediate, since having exactly one solution implies consistency.

$(e)\Rightarrow(a)$: Apply consistency to each of the systems $A\mathbf{x}=\mathbf{e}_1,\dots,A\mathbf{x}=\mathbf{e}_n$ (where $\mathbf{e}_i$ is the $i$th standard unit column). Taking solutions $\mathbf{x}_1,\dots,\mathbf{x}_n$ as columns of a matrix $C=[\mathbf{x}_1\mid\cdots\mid\mathbf{x}_n]$, then $AC=[A\mathbf{x}_1\mid\cdots\mid A\mathbf{x}_n]=[\mathbf{e}_1\mid\cdots\mid\mathbf{e}_n]=I$. By Theorem 1.6.3(b), $C=A^{-1}$, so $A$ is invertible. $\blacksquare$

We know invertible factors give an invertible product; the converse also holds.

\begin{tcolorbox}[colback=red!4, colframe=red!60!black, title=Theorem 1.6.5, boxrule=0.5mm, arc=3mm]
Let $A,B$ be square of the same size. If $AB$ is invertible, then $A$ and $B$ are both invertible.
\end{tcolorbox}
\textbf{Proof} First, $B$ is invertible: if $B\mathbf{x}_0=\mathbf{0}$ for some $\mathbf{x}_0$, then $(AB)\mathbf{x}_0=A(B\mathbf{x}_0)=A\mathbf{0}=\mathbf{0}$; since $AB$ is invertible, Theorem 1.6.4 forces $\mathbf{x}_0=\mathbf{0}$, so $B\mathbf{x}=\mathbf{0}$ has only the trivial solution and $B$ is invertible. Then $A=A(BB^{-1})=(AB)B^{-1}$ is a product of invertible matrices, hence invertible. $\blacksquare$

A fundamental question recurs throughout the text:

\begin{tcolorbox}[colback=blue!5, colframe=blue!60!black, title=A Fundamental Problem, boxrule=0.5mm, arc=3mm]
For a fixed $m\times n$ matrix $A$, find all $m\times1$ matrices $\mathbf{b}$ for which $A\mathbf{x}=\mathbf{b}$ is consistent.
\end{tcolorbox}
If $A$ is invertible, Theorem 1.6.2 fully answers this: every $\mathbf{b}$ works. Otherwise (non-square, or square but singular), $\mathbf{b}$ must usually satisfy some condition, found as in the next examples.

\subsection*{EXAMPLE 3 | Determining Consistency by Elimination}
What must $b_1,b_2,b_3$ satisfy for $x_1+x_2+2x_3=b_1,\ x_1+x_3=b_2,\ 2x_1+x_2+3x_3=b_3$ to be consistent?

\textbf{SOLUTION:} Row-reducing the augmented matrix $\begin{bmatrix}1&1&2&b_1\\1&0&1&b_2\\2&1&3&b_3\end{bmatrix}$ (add $-1\times$row 1 to row 2 and $-2\times$row 1 to row 3; multiply row 2 by $-1$; add row 2 to row 3) yields
\[
\begin{bmatrix}1&1&2&b_1\\0&1&1&b_1-b_2\\0&0&0&b_3-b_2-b_1\end{bmatrix}
\]
so the system is consistent iff $b_3-b_2-b_1=0$, i.e., $b_3=b_1+b_2$ -- that is, iff $\mathbf{b}=\begin{bmatrix}b_1\\b_2\\b_1+b_2\end{bmatrix}$ for arbitrary $b_1,b_2$.

\subsection*{EXAMPLE 4 | Determining Consistency by Elimination}
What must $b_1,b_2,b_3$ satisfy for $x_1+2x_2+3x_3=b_1,\ 2x_1+5x_2+3x_3=b_2,\ x_1+8x_3=b_3$ to be consistent?

\textbf{SOLUTION:} Reducing $\begin{bmatrix}1&2&3&b_1\\2&5&3&b_2\\1&0&8&b_3\end{bmatrix}$ to reduced row echelon form gives
\[
\begin{bmatrix}1&0&0&-40b_1+16b_2+9b_3\\0&1&0&13b_1-5b_2-3b_3\\0&0&1&5b_1-2b_2-b_3\end{bmatrix} \tag{2}
\]
Here $b_1,b_2,b_3$ are unrestricted -- the system is always consistent, with unique solution
\[
x_1=-40b_1+16b_2+9b_3,\quad x_2=13b_1-5b_2-3b_3,\quad x_3=5b_1-2b_2-b_3 \tag{3}
\]
for all $b_1,b_2,b_3$. (This coefficient matrix is precisely the invertible $A$ of Example 1 -- consistent with Theorem 1.6.4(e).)

\section*{Exercise Set 1.6}
\begin{enumerate}
    \item[1.] Solve by inverting the coefficient matrix and using Theorem 1.6.2.\\
    $x_1+x_2=2,\quad 5x_1+6x_2=9$

    \item[3.] $x_1+3x_2+x_3=4,\quad 2x_1+2x_2+x_3=-1,\quad 2x_1+3x_2+x_3=3$

    \item[6.] $-x-2y-3z=0,\quad w+x+4y+4z=7,\quad w+3x+7y+9z=4,\quad -w-2x-4y-6z=6$

    \item[7.] Solve for $x_1,x_2$ in terms of $b_1,b_2$: $3x_1+5x_2=b_1,\quad x_1+2x_2=b_2$

    \item[9.] Solve the linear systems together by reducing an appropriate augmented matrix to reduced row echelon form.\\
    $x_1-5x_2=b_1,\quad 3x_1+2x_2=b_2$\\
    (i) $b_1=1,b_2=4$\qquad (ii) $b_1=-2,b_2=5$

    \item[13.] Determine conditions on the $b_i$'s, if any, for the system to be consistent.\\
    $x_1+3x_2=b_1,\quad -2x_1+x_2=b_2$

    \item[15.] $x_1-2x_2+5x_3=b_1,\quad 4x_1-5x_2+8x_3=b_2,\quad -3x_1+3x_2-3x_3=b_3$

    \item[18.] For $A=\begin{bmatrix}2&1&2\\2&2&-2\\3&1&1\end{bmatrix}$, $\mathbf{x}=\begin{bmatrix}x_1\\x_2\\x_3\end{bmatrix}$:\\
    (a) Show that $A\mathbf{x}=\mathbf{x}$ can be rewritten as $(A-I)\mathbf{x}=\mathbf{0}$, and use this to solve $A\mathbf{x}=\mathbf{x}$.\\
    (b) Solve $A\mathbf{x}=4\mathbf{x}$.

    \item[19.] Solve the matrix equation for $X$.
    \[
    \begin{bmatrix}1&-1&1\\2&3&0\\0&2&-1\end{bmatrix}X=\begin{bmatrix}2&-1&5&7&8\\4&0&-3&0&1\\3&5&-7&2&1\end{bmatrix}
    \]

    \item[21.] \textbf{(Working with Proofs)} Let $A\mathbf{x}=\mathbf{0}$ be a homogeneous system of $n$ equations in $n$ unknowns with only the trivial solution. Prove that for any positive integer $k$, $A^k\mathbf{x}=\mathbf{0}$ also has only the trivial solution.

    \item[23.] \textbf{(Working with Proofs)} Let $A\mathbf{x}=\mathbf{b}$ be any consistent system, and let $\mathbf{x}_1$ be a fixed solution. Prove that every solution can be written $\mathbf{x}=\mathbf{x}_1+\mathbf{x}_0$ where $\mathbf{x}_0$ solves $A\mathbf{x}=\mathbf{0}$, and that every matrix of this form is a solution.

    \item[24.] \textbf{(Working with Proofs)} Use part (a) of Theorem 1.6.3 to prove part (b).
\end{enumerate}

\end{document}
